\documentclass[aspectratio=169]{beamer}

\usetheme{Madrid}
\usecolortheme{default}

\usepackage{graphicx}
\usepackage{booktabs}
\usepackage{amsmath}
\usepackage{adjustbox}
\usepackage{natbib}
\bibliographystyle{chicago}

% Title context
\title[Rural Population Change]{Rural Population Change in Europe}
\subtitle{The Role of Policy, Economic Structure, and Space}
\author[W\"ogerer, Golic, Krisztin]{Michael W\"ogerer \and Dejana Golic \and Tamas Krisztin}
\institute[IIASA]{International Institute for Applied Systems Analysis (IIASA)}
\date{\today}

\begin{document}

% Title Slide
\begin{frame}
    \titlepage
\end{frame}

% Part 1: Introduction and Literature
\begin{frame}{Introduction: The Rural Policy Challenge}
    \textbf{Background \& Motivation:}
    \begin{itemize}
        \item Rural depopulation persists despite decades of place-based EU interventions, contributing to a widening ``geography of discontent'' \citep{RodriguezPose2018, McCann2019}. \pause
        \item Drivers of the rural exodus are well-documented: mechanization lowering primary sector labor demand \citep{Collantes2011} and urban agglomeration economies \citep{Glaeser2009, lichter2025depopulation}. \pause
    \end{itemize}
    \vspace{0.3cm}
    \textbf{Research Gap:}
    \begin{itemize}
        \item The demographic return on EU public investment remains unclear. Aggregating distinct funding streams masks heterogeneous regional outcomes \citep{RodriguezPose2004}. \pause
        \item It remains empirically untested whether the CAP ``Green Deal'' pivot generates co-benefits for population retention or inadvertently crowds out economic activity.
    \end{itemize}
\end{frame}

% Part 2: Data and Stylized Facts
\begin{frame}{Defining the Dependent Variable}
    \textbf{Outcome Measure: Net Migration Rate (2011--2021)}
    \vspace{0.3cm}
    \begin{itemize}
        \item Direct migration flow data is often incomplete at the granular (NUTS~3) level.\pause
        \item We therefore derive net migration using the \textbf{residual approach}:\pause
        \begin{itemize}
            \item Calculate total population change over the decade.
            \item Subtract \textit{natural} population change (cumulative births minus deaths).
        \end{itemize}\pause
        \item The final variable is expressed as a rate relative to the 2011 baseline population.
        \item \textit{Advantage:} This strictly isolates mobility-driven shifts from fertility/mortality trends.
    \end{itemize}
\end{frame}

\begin{frame}{Stylized Facts: Spatial Clustering of Population Change}
    \centering
    \includegraphics[height=0.85\textheight]{figures/combined_growth_corr.png}
\end{frame}

\begin{frame}{Stylized Facts: Depopulation by Regional Typology}
    \begin{columns}
        \begin{column}{0.5\textwidth}
            \centering
            \includegraphics[width=\textwidth]{figures/relative_density_corr.png}
        \end{column}
        \begin{column}{0.5\textwidth}
            \textbf{Key Observations:}
            \begin{itemize}
                \item Depopulation is predominantly a rural phenomenon, with the rural density distribution (green) shifted substantially to the left.
                \item Population stability and growth are concentrated within urban and intermediate regions.
                \item This strong spatial clustering inherently motivates the subsequent use of spatial econometric modeling.
            \end{itemize}
        \end{column}
    \end{columns}
\end{frame}

% Part 3: Policy Classification & Estimation
\begin{frame}{Data Overview: Regional Drivers of Population Change}
    To isolate policy effects, we control for baseline regional endowments (circa 2011 / 2009--2013 average) to mitigate simultaneity bias.
    \vspace{0.3cm}
    \begin{columns}[T]
        \begin{column}{0.5\textwidth}
            \textbf{Structural Economic Drivers:}
            \begin{itemize}
                \item GDP per capita (constant prices).
                \item Employment (jobs per capita).
                \item GVA Shares (Agriculture, Mining/Utilities) to capture structural transformation.
            \end{itemize}
        \end{column}
        \begin{column}{0.5\textwidth}
            \textbf{Amenity \& Geography:}
            \begin{itemize}
                \item Accessibility Index (inverted drive times to education/health).
                \item Population Density (and squared term).
                \item Fixed effects: Coastal, Capital, and Metro typologies.
            \end{itemize}
        \end{column}
    \end{columns}
\end{frame}

\begin{frame}{Breaking Down the EU Budget}
    To understand policy impacts, we must unbundle aggregate funds into their functional mechanisms:
    \vspace{0.5cm}
    \begin{columns}
        \begin{column}{0.5\textwidth}
            \textbf{CAP (Common Agricultural Policy)}
            \begin{itemize}
                \item \textit{Pillar I:} Direct income support (Welfare floor).
                \item \textit{Pillar II (Development):} Modernization, diversification.
                \item \textit{Pillar II (Environment):} Agri-environmental schemes (Natura 2000).
            \end{itemize}
        \end{column}
        \begin{column}{0.5\textwidth}
            \textbf{ESIF (Structural Funds)}
            \begin{itemize}
                \item \textit{Construction:} Hard infrastructure (TEN-T, transport).
                \item \textit{Environment:} Risk prevention, brownfield redevelopment.
            \end{itemize}
        \end{column}
    \end{columns}
    \vspace{0.5cm}
    \textit{Note: All policies measured as intensity (Log Amount / GDP per capita) to capture relative economic importance.}
\end{frame}

% Part 3: The Methodological Engine
\begin{frame}{The Spatial Autoregressive (SAR) Approach}
    \textbf{Why not OLS?} Demography is not spatially contained. A region imports/exports population stability from neighbors.
    
    \vspace{0.3cm}
    \textbf{Structural Model:}
    \begin{equation}
    \mathbf{y} = \rho \mathbf{W} \mathbf{y} + \mathbf{P} \boldsymbol{\gamma} + \mathbf{X} \boldsymbol{\beta} + \boldsymbol{\varepsilon}, \quad \boldsymbol{\varepsilon} \sim \mathcal{N}(0, \sigma^2 \mathbf{I}_N)
    \end{equation}

    {\footnotesize
    \begin{itemize}
        \item \textbf{$\mathbf{y}$}: Vector of average annual net migration rates (2011--2021).\pause
        \item \textbf{$\rho$}: Spatial autoregressive parameter capturing cross-regional spillovers.
        \item \textbf{$\mathbf{W}$}: Row-stochastic $k$-nearest-neighbor spatial weights matrix ($k=10$).\pause
        \item \textbf{$\mathbf{P}$ \text{and} $\boldsymbol{\gamma}$}: Matrix of EU policy exposure intensities (CAP, ESIF) and their estimated effects.\pause
        \item \textbf{$\mathbf{X}$ \text{and} $\boldsymbol{\beta}$}: Matrix of baseline economic covariates and their estimated effects.\pause
        \item \textbf{$\boldsymbol{\varepsilon}$}: Normally distributed error term.
    \end{itemize}
    
    \vspace{0.3cm}
    \textbf{Note:} Estimation is carried out via \textbf{Bayesian MCMC} to naturally accommodate spatial parameter and coefficient uncertainty simultaneously.
    }
\end{frame}

\begin{frame}{Dealing with Endogeneity}
    \textbf{The problem:} Struggling regions receive more compensatory funding (simultaneity bias).\pause
    \vspace{0.5cm}
    
    \textbf{Our Long-Difference Design:}
    \begin{itemize}
        \item Regress decadal population change (2011--2021)...
        \item ...on policy exposure and covariates measured at the \textbf{baseline} (2011 / 2009--2013 average).
        \item Ensures explanatory variables are predetermined relative to the subsequent demographic trajectory.
    \end{itemize}
\end{frame}

% Part 4: Core Findings & Regression Results
\begin{frame}{Regression Strategy: Three Nested Models}
    To isolate the differential impacts of policy across regional typologies, we estimate three nested specifications:\pause
    \begin{itemize}
        \item \textbf{Model (i) - Base Model:} Controls for structural, demographic, and geographic drivers to establish the baseline determinants of population change.\pause
        \item \textbf{Model (ii) - Policy Average Effects:} Introduces aggregated CAP and ESIF funding intensities to measure average EU-wide policy efficacy.\pause
        \item \textbf{Model (iii) - Typology Interactions:} Interacts policy intensities with regional typologies (Urban, Intermediate, Rural) to uncover spatial heterogeneity in policy impact.
    \end{itemize}
\end{frame}

\begin{frame}{Results: Structural Covariates \& Spatial Spillovers}
    % Shows the variables that remain robust across all three specifications
    \small
    \input{bigsar_structural.tex}
\end{frame}

\begin{frame}[shrink=20]{Results: Policy Drivers \& Regional Interactions}
    % Shows the policy blocks
    \footnotesize
    \input{bigsar_policy.tex}
\end{frame}

\begin{frame}{Core Finding 1: Structural Destiny \& Spillovers}
    % SPEAKING NOTE: "Agriculture GVA" is not monolithic. Intensive agriculture (e.g., horticulture) might attract labor, 
    % while mechanized arable farming drives the "demographic tax" through labor shedding.
    \begin{itemize}
        \item \textbf{The Demographic Tax of Agriculture:} A 10 percentage point increase in agricultural GVA share is associated with a 3.1 percentage point reduction in population growth. \pause
        \item \textbf{The Metropolitan Premium:} Capital and metro regions enjoy huge baseline advantages.
        % SPEAKING NOTE: Explicitly mention that because we control for accessibility and proximity to metros, 
        % our policy interaction terms in the next slide isolate the effect on 'deep' rural areas, not just commuting zones.
        \pause
        \item \textbf{Massive Spatial Spillovers:} $\rho \approx 0.67$. Two-thirds of variation in local growth rates is spatially clustered.
    \end{itemize}
\end{frame}

\begin{frame}{Core Finding 2: The ``Green-Demographic'' Trade-off}
    Looking specifically at the interaction model for Rural regions:
    \vspace{0.3cm}
    % SPEAKING NOTE: Explain the "Lifeline" mechanism. Even though Pillar I goes to farmers (often large landowners), 
    % it subsidizes local rural consumption and sustains local services, acting as a regional macroeconomic stabilizer.
    \begin{block}{Pillar I: The Lifeline}
        Direct income support (Pillar I) acts as a vital demographic stabilizer for rural areas, vastly outperforming ESIF hard infrastructure.
    \end{block}
    \vspace{0.2cm}\pause
    % SPEAKING NOTE: Address the "Green Trade-off". Explain that Agri-environmental schemes often require 
    % de-intensification or land abandonment (rewilding). This reduces the carrying capacity for traditional 
    % agricultural labor faster than alternative 'green' jobs can replace them.
    \begin{block}{Pillar II Environment: The Provocation}
        \textbf{Significant negative association} between Agri-environmental schemes and rural population retention. 
    \end{block}
    \vspace{0.2cm}
    \textit{Insight:} Maximizing ecological utility (re-wilding, extensive farming) may inadvertently hollow out rural labor demand, accelerating the rural exodus.
\end{frame}

% Part 5: Conclusion
\begin{frame}{Conclusion and Policy Implications}
    \textbf{The Policy Trilemma:}
    \begin{itemize}
        \item It is exceedingly difficult to simultaneously achieve:
        \begin{enumerate}
            \item Rural population retention
            \item Agricultural modernization
            \item Aggressive environmental conservation (Green Deal)
        \end{enumerate}
    \end{itemize}
    \vspace{0.5cm}
    \textbf{Takeaway:}
    As the EU pivots towards a ``green'' CAP, stripping direct income support (Pillar I) to fund environmental schemes (Pillar II) without providing alternative non-agricultural employment will likely accelerate rural depopulation.
\end{frame}

\begin{frame}{Limitations and Future Research}
    \begin{itemize}
        % SPEAKING NOTE: Address Identification. Framing it as conditional spatial associations vs. a perfect experiment.
        \item \textbf{Identification Strategy:} While the baseline-lagged, long-difference approach mitigates simultaneity, it yields conditional spatial associations rather than a fully structural model of migration.\pause
        % SPEAKING NOTE: Address Migration Quality vs Quantity. Acknowledge that net migration captures numbers but not age distribution or human capital quality.
        \item \textbf{Migration Quality vs. Quantity:} Our dependent variable ($growth\_corr$) is aggregate net migration. It captures raw numbers but does not untangle the demographic composition (e.g., retaining retirees vs. losing young, skilled labor).\pause
        \item \textbf{Future Directions:} Structural modeling is needed to determine if ``Green'' CAP payments differentially affect the retention of youth versus older populations.
    \end{itemize}
\end{frame}

\begin{frame}
    \centering
    \Huge
    Thank You! \\
    \vspace{0.5cm}
    \Large
    Questions?
\end{frame}

\begin{frame}[allowframebreaks]{References}
\scriptsize
\begin{thebibliography}{99}

\bibitem[Camaioni et al.(2016)]{Camaioni2016}
Camaioni, B., Esposti, R., Lobianco, A., Pagliacci, F., and Sotte, F. (2016).
The CAP Vademecum: Abstracting the European agricultural policy.
\textit{Agecon Search}.

\bibitem[Collantes and Pinilla(2011)]{Collantes2011}
Collantes, F., and Pinilla, V. (2011).
\textit{Peaceful surrender: The depopulation of rural Spain in the twentieth century}.
Cambridge Scholars Publishing.

\bibitem[Glaeser and Gottlieb(2009)]{Glaeser2009}
Glaeser, E. L., and Gottlieb, J. D. (2009).
The wealth of cities: Agglomeration economies and spatial equilibrium in the United States.
\textit{Journal of Economic Literature}, 47(4), 983--1028.

\bibitem[Lichter and Johnson(2025)]{lichter2025depopulation}
Lichter, D. T., and Johnson, K. M. (2025).
Depopulation, Deaths, Diversity, and Deprivation: The 4Ds of Rural Population Change.
\textit{RSF: The Russell Sage Foundation Journal of the Social Sciences}, 11(2), 88--114.

\bibitem[McCann(2020)]{McCann2019}
McCann, P. (2020).
Perceptions of regional inequality and the geography of discontent: insights from the UK.
\textit{Regional Studies}, 54(2), 256--267.

\bibitem[Rodr\'{\i}guez-Pose(2018)]{RodriguezPose2018}
Rodr\'{\i}guez-Pose, A. (2018).
The revenge of the places that don't matter (and what to do about it).
\textit{Cambridge Journal of Regions, Economy and Society}, 11(1), 189--209.

\bibitem[Rodr\'{\i}guez-Pose and Fratesi(2004)]{RodriguezPose2004}
Rodr\'{\i}guez-Pose, A., and Fratesi, U. (2004).
Between development and social policies: the impact of European Structural Funds in Objective 1 regions.
\textit{Regional Studies}, 38(1), 97--113.

\end{thebibliography}
\end{frame}

\end{document}
